\section{Indledning}

\textbf{Terraen\_model} er en QGIS Processing-model der automatisk opbygger en
hydrologisk konditioneret terrænmodel til brug i vådområdemodellering. Modellen
henter det danske højdepunktsnet (DHM Terræn) direkte fra Kortforsyningen,
brænder drænlinjer (DHMLinje) ind i rastermodellen for at sikre korrekte
afstrømningsveje, og udleder dernæst et kanalsystem til brug i videre
oplandsbereegninger.

Resultatet er tre outputfiler der gemmes automatisk i projektmappens
\texttt{Outputfiler\_<omraade>/}-undermappe:

\begin{itemize}
  \item \texttt{Hoejdemodel.sdat} — Den konditionerede terrænmodel (SAGA-format)
  \item \texttt{Channel\_Networks.gpkg} — Udledt kanalsystem
  \item \texttt{DHMLinje\_Klippet.gpkg} — Klippede DHMLinje-linjer (én pr. parallel gruppe)
\end{itemize}

\textbf{Overordnet arbejdsgang:}

\begin{enumerate}
  \item Download DHM Terræn for projektområdet (inkl.\ 3 km buffer)
  \item Hent og klip DHMLinje til samme buffer
  \item Udvælg geometrisk centrale linjer der, hvor parallelle linjer overlapper
  \item Drappe DHM's z-værdier på de udvalgte linjer
  \item Beregn interpolerede kanalbundshøjder (\texttt{z\_interp}) langs linjerne
  \item Brænd en V-formet dale ind i DHM med kosinusbaseret overgangszone (0–30 m)
  \item Erstat kanalbunden direkte (4 m zone) med de interpolerede z-værdier
  \item Fyld kunstige sænkninger (Fill Sinks)
  \item Udled kanalsystem og drænoplande
  \item Gem output
\end{enumerate}


\section{Inputparametre}

\begin{tabular}{lp{10cm}}
  \hline
  \textbf{Parameter} & \textbf{Beskrivelse} \\
  \hline
  \texttt{projektomraade} & Et polygon-lag der definerer projektområdets udstrækning.
    Modellen opretter en 3 km buffer og begrænser alle beregninger til dette udsnit. \\[4pt]
  \texttt{token} & Adgangstoken til Kortforsyningen (Geodatastyrelsen).
    Bruges til at hente DHM Terræn-data via WCS. \\
  \hline
\end{tabular}

\medskip
DHMLinje hentes automatisk fra mappehierarkiet og kræver ingen manuel
parameterindstilling (se afsnit~\ref{sec:indput_dhmlinje}).


\section{Modelkæden — trin for trin}

\subsection{Trin 1 — Buffer (6000 m)}

\begin{tabular}{ll}
  \textit{Algoritme:} & \texttt{native:buffer} \\
  \textit{Input:}     & \texttt{projektomraade} \\
  \textit{Distance:}  & 6000 m \\
\end{tabular}

\medskip
Projektområdet udvides med 6 km i alle retninger. Den store buffer sikrer, at
DHM-downloaden og de efterfølgende hydrologiske beregninger ikke rammes af
kanteffekter. Fill-Sinks-algoritmen er især følsom over for kortede
terrænmodeller, idet kunstige ``vægge'' langs rasterets kant kan tvinge
afstrømning i forkerte retninger. En bred buffer er særligt vigtig for
store, flade vådområder med lange afstrømningsveje.

\subsection{Trin 2 — Hent DHM Terræn (WCS)}

\begin{tabular}{ll}
  \textit{Algoritme:}  & \texttt{script:hent\_dhm\_wcs} \\
  \textit{Service:}    & \texttt{api.dataforsyningen.dk/dhm\_wcs\_DAF} \\
  \textit{Coverage:}   & \texttt{dhm\_terraen} \\
  \textit{Opløsning:}  & 2 m/pixel (standard) \\
  \textit{Projektion:} & EPSG:25832 (UTM zone 32N) \\
\end{tabular}

\medskip
DHM Terræn er det danske LiDAR-baserede digitale terrænmodel med en original
opløsning på ca.\ 0,4 m/pixel. Modellen downloader det relevante udsnit
svarende til 3 km-bufferen. Opløsningen er som standard sat til 2 m/pixel for
at balancere detaljegrad og procestid; den kan nedjusteres til 0,4 m for høj
præcision eller op til 10 m for oversigtsformål.

GDAL's WCS-driver bruges til at åbne servicen via en XML-konfiguration gemt i
GDAL's virtuelle filsystem (\texttt{/vsimem/}). Derefter udklippes og resamples
det ønskede udsnit med \texttt{gdal.Translate}.

\subsection{Trin 3 — Indput DHMLinje (automatisk)}

\begin{tabular}{ll}
  \textit{Algoritme:} & \texttt{script:indput\_dhmlinje} \\
  \textit{Kilde:}     & \texttt{Grunddata/DHMLinje/DHMLinje.shp} \\
\end{tabular}

\medskip
DHMLinje er et linjekortlag der repræsenterer vandløb og drænkanaler.
I stedet for at brugeren manuelt skal pege på filen, finder scriptet automatisk
\texttt{DHMLinje.shp} ved at søge opad og nedad i mappehierarkiet fra enten en
brugerkonfigureret rodmappe (\texttt{QgsSettings}) eller QGIS-projektets
hjemmemappe. Se afsnit~\ref{sec:indput_dhmlinje} for detaljer.

\subsection{Trin 4 — Klip DHMLinje til buffer}

\begin{tabular}{ll}
  \textit{Algoritme:} & \texttt{native:clip} \\
  \textit{Input:}     & DHMLinje (fra trin 3) \\
  \textit{Overlay:}   & 3000 m buffer (fra trin 1) \\
\end{tabular}

\medskip
Begrænser DHMLinje-laget til det relevante undersøgelsesområde og reducerer
efterfølgende procestid i trin 5–14.

\subsection{Trin 5 — Udvælg midterste linjer}

\begin{tabular}{ll}
  \textit{Algoritme:}   & \texttt{script:udpeg\_midterste\_linjer} \\
  \textit{Søgeafstand:} & 8 m \\
\end{tabular}

\medskip
DHMLinje-datasættet indeholder mange steder to eller flere parallelle linjer der
repræsenterer det samme vandløb. Hvis alle linjer bruges i den efterfølgende
V-profil-korrektion, opstår der sammenstødende bufferzoner og inkonsistente
z-værdier.

Scriptet grupperer linjer der ligger inden for 8 m af hinanden (buffer +
dissolve + singleparts) og beholder for hver gruppe kun den linje, hvis
tyngdepunkt er geometrisk tættest på gruppens samlede tyngdepunkt. Se
afsnit~\ref{sec:midterste_linjer} for detaljer.

\subsection{Trin 6 — Drappe Z fra DHM}

\begin{tabular}{ll}
  \textit{Algoritme:} & \texttt{native:setzfromraster} \\
  \textit{Input:}     & Midterste linjer (fra trin 5) \\
  \textit{Raster:}    & DHM Terræn (fra trin 2) \\
  \textit{Band:}      & 1 \\
  \textit{NODATA:}    & $-9999$ \\
\end{tabular}

\medskip
Tildeler en z-koordinat til hvert knudepunkt på de udvalgte linjer ved at slå
DHM-højden op i den tilsvarende rastercelle. Outputlaget er 3D-linjer, hvor
hvert vertex har en z-værdi der afspejler det aktuelle terræn.

Drapping-trinnet er nødvendigt fordi \texttt{native:pointsalonglines} (trin 8)
udlæser z-koordinaten fra linjernes geometri som \texttt{\$z}-attributten.
Uden dette trin ville alle z-værdier være 0 eller udefinerede, og
z\_interp-beregningen i trin 9 ville give NULL for alle punkter.

\subsection{Trin 7 — Beregn linjelængde}

\begin{tabular}{ll}
  \textit{Algoritme:} & \texttt{native:fieldcalculator} \\
  \textit{Nyt felt:}  & \texttt{line\_length} \\
  \textit{Udtryk:}    & \texttt{\$length} \\
\end{tabular}

\medskip
Tilføjer linjelængden som et attributfelt til hver DHMLinje-feature. Feltet
bruges i trin 8 til at generere samplingpunkter med korrekt afstandsangivelse.

\subsection{Trin 8 — Punkter langs linjer (2 m)}

\begin{tabular}{ll}
  \textit{Algoritme:} & \texttt{native:pointsalonglines} \\
  \textit{Interval:}  & 2 m \\
\end{tabular}

\medskip
Genererer samplingpunkter langs hver DHMLinje-feature med 2 m mellemrum.
Hvert punkt arver linjens 3D-geometri fra trin 6, og z-koordinaten er
tilgængelig via QGIS-udtrykket \texttt{\$z}. Algoritmen tilføjer desuden
attributten \texttt{distance} (afstand fra linjens startpunkt).

\subsection{Trin 9 — Beregn z\_interp}

\begin{tabular}{ll}
  \textit{Algoritme:} & \texttt{native:fieldcalculator} \\
  \textit{Nyt felt:}  & \texttt{z\_interp} \\
  \textit{Udtryk:}    & \texttt{"START\_Z" + ("END\_Z" - "START\_Z") * ("distance" / "line\_length")} \\
\end{tabular}

\medskip
Forsøger at beregne en lineær kanalbundshøjde der interpoleres fra linjens
startpunkt (\texttt{START\_Z}) til slutpunkt (\texttt{END\_Z}) proportionalt
med afstanden \texttt{distance}. Hensigten er at bruge DHM-terrænet i endepunkterne
som reference frem for de rå, potentielt støjbehæftede punkthøjder langs linjen.

\medskip
\textbf{Praktisk begrænsning:} \texttt{START\_Z} og \texttt{END\_Z} er ikke
standardattributter fra \texttt{native:pointsalonglines}. Hvis inputlinjerne ikke
allerede har disse felter, returnerer feltberegneren NULL for alle samplingpunkter.
I praksis bidrager z\_interp-rasteret (trin 11) derfor ikke til den endelige
terrænmodel — terrænkonditioneringen baseres udelukkende på V-profil-sænkningen
fra trin 14. Se afsnit~\ref{sec:begr} (kendte begrænsninger) for perspektiv.

\subsection{Trin 10 — Buffer (4 m) om z\_interp-punkter}

\begin{tabular}{ll}
  \textit{Algoritme:} & \texttt{native:buffer} \\
  \textit{Distance:}  & 4 m \\
\end{tabular}

\medskip
Opbygger en 4 m bred zone langs kanalen. Inden for denne zone erstattes DHM
direkte med de interpolerede z\_interp-værdier (trin 11 og 15). Zone-bredden
svarer til kanalbundens umiddelbare omgivelser og sikrer, at den brændte kanal
er geometrisk sammenhængende i det endelige raster.

\subsection{Trin 11 — Rasterize z\_interp}

\begin{tabular}{ll}
  \textit{Algoritme:} & \texttt{gdal:rasterize} \\
  \textit{Burn felt:} & \texttt{z\_interp} \\
  \textit{NODATA:}    & $-9999$ \\
\end{tabular}

\medskip
Konverterer z\_interp-punkternes 4 m bufferzoner til et raster med samme
opløsning og udstrækning som DHM. Celler inden for bufferzonerne får
z\_interp-værdien; alle øvrige celler sættes til NODATA ($-9999$).

\subsection{Trin 12 — Rasterize DHMLinje (binær)}

\begin{tabular}{ll}
  \textit{Algoritme:} & \texttt{gdal:rasterize} \\
  \textit{BURN:}      & 1 \\
  \textit{INIT:}      & 0 \\
\end{tabular}

\medskip
Opretter et binært raster (0/1) der angiver, hvilke celler der ligger direkte
på en DHMLinje. Outputtet bruges i trin 13 som kilde til afstandsberegningen.

\subsection{Trin 13 — Nærhedsraster (proximity)}

\begin{tabular}{ll}
  \textit{Algoritme:}     & \texttt{gdal:proximity} \\
  \textit{MAX\_DISTANCE:} & 25 m \\
  \textit{NODATA:}        & 26 \\
\end{tabular}

\medskip
Beregner for hver rastercelle den korteste afstand til nærmeste DHMLinje-celle.
Celler mere end 25 m væk tildeles NODATA-værdien 26, hvilket sikrer, at
V-profilkorrektionen ikke påvirker terrænnet uden for en 25 m-zone.
Zonebredden er afstemt med den hydrologiske bremseeffekt: kanalen giver kun
meningsfuld drænpåvirkning inden for ca.\ 25 m fra vandløbslinjen.

\subsection{Trin 14 — Rasterberegner: kosinusbaseret V-profil}

\begin{tabular}{ll}
  \textit{Algoritme:} & \texttt{gdal:rastercalculator} \\
  \textit{Raster A:}  & DHM (fra trin 2) \\
  \textit{Raster B:}  & Nærhedsraster (fra trin 13) \\
\end{tabular}

\medskip
Formlen:
\begin{verbatim}
A - numpy.maximum(0.0,
      1.0 + numpy.cos(numpy.pi * numpy.minimum(B, 25.0) / 25.0))
\end{verbatim}

Dette er kernen i terrænkonditioneringen. Den skaber en blød V-formet
sænkning i DHM, centreret om DHMLinje, med:

\begin{itemize}
  \item \textbf{Maksimal korrektion ved afstand = 0:}
        $1 + \cos(0) = 2$ meter sænkning
  \item \textbf{Nul-korrektion ved afstand = 25 m:}
        $1 + \cos(\pi) = 0$ meter ændring
  \item \textbf{Glat overgang:} Kosinusfunktionen giver en S-kurveformet
        aftrapning, som undgår de skarpe kanter der opstår ved lineær aftrapning.
  \item \textbf{Beskyttelse mod opløftning:}
        \texttt{numpy.maximum(0, \ldots)} sikrer, at formlen aldrig hæver terrænet.
\end{itemize}

Matematisk:
\[
  z'(d) = z(d) - \max\!\left(0,\; 1 + \cos\!\left(\frac{\pi \cdot \min(d,\,25)}{25}\right)\right)
\]
hvor $d$ er afstand til DHMLinje i meter og $z$ er den originale DHM-højde.

\subsection{Trin 15 — Merge: kombiner V-profil og z\_interp}

\begin{tabular}{ll}
  \textit{Algoritme:}        & \texttt{gdal:merge} \\
  \textit{Lag (rækkefølge):} & {[V-profil raster, z\_interp raster]} \\
  \textit{NODATA\_INPUT:}    & $-9999$ \\
\end{tabular}

\medskip
z\_interp-rasteret (trin 11) har NODATA for alle celler uden for 4 m-bufferen,
og V-profil-rasteret (trin 14) dækker hele fladen. Ved at liste z\_interp-rasteret
sidst i stakken \emph{overskriver} dets gyldige celler V-profil-værdierne i
kanalbunden, mens V-profil-rasteret udfylder resten.

Det endelige resultat er et DHM der:
\begin{enumerate}
  \item Er uændret mere end 30 m fra DHMLinje
  \item Er glat sænket i 0–30 m-zonen (kosinusovergang)
  \item Er erstattet med faktiske DHM-kanalbundsværdier i 0–4 m-zonen
\end{enumerate}

\subsection{Trin 16 — Fill Sinks (Wang \& Liu)}

\begin{tabular}{ll}
  \textit{Algoritme:}  & \texttt{native:fillsinkswangliu} \\
  \textit{Input:}      & Merge-raster (trin 15) \\
  \textit{MIN\_SLOPE:} & 0,01 \\
\end{tabular}

\medskip
SAGA's Wang \& Liu-algoritme fjerner kunstige sænkninger i terrænmodellen som
hindrer sammenhængende afstrømningsveje. Et minimumsgradient på 0,01 sikrer at
alle overflader har en svag hældning og derved afvander til en kant frem for at
samle i et lokalt lavpunkt.

Fill Sinks er nødvendigt selv efter V-profilkorrektionen, fordi lokale
LiDAR-artefakter (broer, bygninger, vegetationssupport) kan skabe lukkede
sænkninger, og DHM-rasterets kantceller skal have klare afstrømningsveje.

\subsection{Trin 17 — Kanalsystem og drænoplande}

\begin{tabular}{ll}
  \textit{Algoritme:} & \texttt{sagang:channelnetworkanddrainagebasins} \\
  \textit{THRESHOLD:} & 5 \\
\end{tabular}

\medskip
SAGA's \textit{Channel Network and Drainage Basins}-modul beregner akkumuleret
afstrømningsareal og udleder kanaler der, hvor den akkumulerede afstrømning
overstiger tærskelværdien 5. Tærskelværdien styrer kanalnetværkets densitet:
lavere værdier giver tættere net, højere værdier giver grovere net.

\subsection{Trin 18 — Gem output}

\begin{tabular}{ll}
  \textit{Algoritme:} & \texttt{script:terraen\_model\_output} \\
\end{tabular}

\medskip
Gemmer de tre resultater i projektmappens \texttt{Outputfiler\_<omraade>/}:
\begin{itemize}
  \item \texttt{Hoejdemodel.sdat} — Det konditionerede raster konverteret til
        SAGA-format med \texttt{gdal:translate}
  \item \texttt{Channel\_Networks.gpkg} — Kanalsystemet som GeoPackage
  \item \texttt{DHMLinje\_Klippet.gpkg} — De udvalgte DHMLinje-linjer (én pr.
        parallel gruppe) med drappede z-værdier, som GeoPackage
\end{itemize}
Scriptet håndterer Windows-fillåse ved at fjerne indlæste lag, køre Python's
garbage collector og afvente GDAL's OGR-forbindelseslukninger.


\section{Hjælpescripts}

Alle scripts ligger i mappen \texttt{Scripts/} og registreres automatisk i
QGIS Processing ved opstart via \texttt{modelmappe\_interface.py}.

\subsection{hent\_dhm\_wcs\_script.py}

Downloader DHM Terræn fra Kortforsyningen via GDAL's WCS-driver.

\begin{itemize}
  \item Bygger en WCS-XML-konfiguration i GDAL's virtuelle filsystem
        (\texttt{/vsimem/dhm\_wcs.xml}) med service-URL og coverage-navn
  \item Åbner servicen som et virtuelt GDAL-datasæt
  \item Udklippes og resamples det ønskede udsnit med \texttt{gdal.Translate}
  \item Understøtter opløsninger fra 0,4 m (original LiDAR) til 10 m
\end{itemize}

\textbf{Hvorfor scriptet er nødvendigt:} QGIS's standardalgoritmer kan hente
WMS/WFS, men ikke WCS i rasterformat med brugerdefineret opløsning og automatisk
klipning til et bbox — det kræver GDAL-API direkte.

\subsection{Indput\_DHMLinje.py}
\label{sec:indput_dhmlinje}

Finder og returnerer stien til \texttt{DHMLinje.shp} automatisk uden brugerinput.

Søgelogikken er trelaget:
\begin{enumerate}
  \item Tjekker brugerkonfigureret rodmappe fra \texttt{QgsSettings}
        (\texttt{vaadomraade\_modeller/mappe}), inkl.\ rekursiv søgning nedad
        i op til 4 mappeniveauer
  \item Søger opad fra QGIS-projektets hjemmemappe i op til 6 niveauer
  \item Kaster \texttt{QgsProcessingException} med vejledende besked hvis
        ingen søgning lykkes
\end{enumerate}

\textbf{Hvorfor scriptet er nødvendigt:} En hardkodet sti ville kræve manuel
tilpasning ved hvert projekt. Søgelogikken gør modellen portabel på tværs af
maskiner og projektmappestrukturer, blot \texttt{Grunddata/DHMLinje/} er til
stede et sted i hierarkiet.

\subsection{Udpeg\_Midterste\_Linjer.py}
\label{sec:midterste_linjer}

Udvælger én repræsentativ linje per gruppe af parallelle linjer.

\textbf{Algoritme:}
\begin{enumerate}
  \item Buffer på \texttt{distance} meter (standard: 8 m) om alle inputlinjer
        + dissolve til ét multi-polygon-objekt
  \item \texttt{native:multiparttosingleparts} splitter multi-polygon'en op i
        individuelle polygoner — én per sammenhængende klynge af parallelle linjer
  \item For hver gruppe-polygon beregnes dens centroid
  \item Alle inputlinjer der overlapper polygonen evalueres; den hvis centroid
        er nærmest polygon-centroiden markeres som vinder
  \item Kun vinderlinjerne skrives til outputlaget
\end{enumerate}

\textbf{Resultat ved $n$ parallelle linjer:}
\begin{itemize}
  \item $n$ ulige: den geometrisk midterste linje vælges
  \item $n$ lige: den linje der er nærmest midtpunktet vælges
  \item $n = 1$: linjen passerer uændret igennem
\end{itemize}

\textbf{Hvorfor multiparttosingleparts er nødvendigt:} \texttt{native:buffer}
med \texttt{DISSOLVE: True} returnerer et enkelt multi-polygon-objekt når blot
to buffere overlapper. Uden opsplitning behandles hele DHMLinje-netværket som
én gruppe, og kun én enkelt linje overlever filtreringen.

\textbf{Hvorfor scriptet er nødvendigt:} DHMLinje-datasættet indeholder hyppigt
to (eller flere) parallelle digitaliserede linjer for samme vandløb. Bruges alle
linjer, opstår der overlappende V-profil-zoner der skaber uregelmæssige
fordybninger og rygge i det konditionerede raster.

\subsection{Terraen\_model\_output.py}

Gemmer modelresultaterne til disk og håndterer Windows-specifikke fillåseproblemer.

\begin{itemize}
  \item Finder output-mappen via \texttt{QgsSettings}:
        \texttt{<rod>/<projekt>/<omraade>/Outputfiler\_<omraade>/}
  \item Konverterer rasteret til SAGA-format (\texttt{.sdat}) med
        \texttt{gdal:translate}
  \item Gemmer \texttt{DHMLinje\_Klippet.gpkg} med
        \texttt{QgsVectorFileWriter.writeAsVectorFormatV3}
  \item Håndterer Windows-fillåse på GPKG-filer med fjernelse af lag fra
        projektet, \texttt{gc.collect()}, OGR \texttt{DeleteDataSource} og
        en retry-løkke med korte pauser
\end{itemize}


\section{Outputfiler}

\begin{tabular}{lp{9cm}}
  \hline
  \textbf{Fil} & \textbf{Indhold} \\
  \hline
  \texttt{Hoejdemodel.sdat} &
    Konditioneret og sink-fyldt terrænmodel i SAGA-format. Bruges direkte i
    SAGA-baserede hydrologimodeller og kan indlæses i QGIS som et rasterbånd. \\[4pt]
  \texttt{Channel\_Networks.gpkg} &
    Udledt kanalsystem som vektorlinjer. Inkluderer SAGA's standardattributter:
    flow accumulation, kanalbrede-estimat m.fl. \\[4pt]
  \texttt{DHMLinje\_Klippet.gpkg} &
    De udvalgte DHMLinje-linjer klippet til 3 km-bufferen (én pr. parallel
    gruppe), med drappede z-værdier fra DHM Terræn. Bruges til visuel kontrol
    og kan danne grundlag for videre hydraulisk modellering. \\
  \hline
\end{tabular}


\section{Forudsætninger og opsætning}

\begin{enumerate}
  \item \textbf{Projektmappestruktur:} Interfacet (\texttt{modelmappe\_interface.py})
        skal have registreret rodmappen under \texttt{QgsSettings}-nøglerne
        \texttt{vaadomraade\_modeller/mappe}, \texttt{/projektnavn} og
        \texttt{/projektomraade\_navn}.

  \item \textbf{Grunddata:} Mappen \texttt{Grunddata/DHMLinje/DHMLinje.shp}
        skal være tilgængelig i hierarkiet under rodmappen.

  \item \textbf{Kortforsyningen-token:} En gyldig API-token er påkrævet for
        at downloade DHM Terræn. Tokens udstedes gratis på
        \texttt{dataforsyningen.dk}.

  \item \textbf{QGIS-plugins:} SAGA-algoritmer kræver SAGA installeret og
        registreret i QGIS Processing
        (\textit{Processing} $\to$ \textit{Options} $\to$ \textit{Providers}
        $\to$ \textit{SAGA}).

  \item \textbf{QGIS-projektet} skal være gemt (\textit{Navngiv projekt} i
        interfacet), og \textit{Tilføj projektområde} skal have oprettet
        outputmappen, inden \texttt{Terraen\_model} køres.
\end{enumerate}


\section{Kendte begrænsninger}
\label{sec:begr}

\begin{itemize}
  \item \textbf{z\_interp virker ikke uden START\_Z/END\_Z-felter:}
        Den lineære interpolationsformel i trin 9 forudsætter at inputlaget
        har attributterne \texttt{START\_Z} og \texttt{END\_Z}. Da
        \texttt{native:pointsalonglines} ikke tilføjer disse felter automatisk,
        er \texttt{z\_interp} NULL for alle punkter i praksis. Det medfører at
        trin 11 (rasterize z\_interp) producerer et raster med udelukkende
        NODATA, og at trin 15 (merge) alene baseres på V-profil-sænkningen.
        For fuld funktionalitet bør formlen erstattes med \texttt{\$z}, der
        aflæser Z-koordinaten direkte fra punktets geometri (sat af
        \texttt{setzfromraster} i trin 6).

  \item \textbf{Afbrudte linjesegmenter:} Hvis DHMLinje-segmenter ikke er
        forbundne ende til ende, kan kanalsystemet (trin 17) vælge
        afstrømningsveje der omgår de brændte kanaler. Løsningen er at sikre
        geometrisk sammenhæng i inputdatasættet.

  \item \textbf{Z-værdier fra DHM:} \texttt{z\_interp} afspejler den faktiske
        DHM-højde i hvert samplingpunkt. Ved linjer der krydser vejdæmninger
        eller broer kan z-værdien være forhøjet relativt til den egentlige
        kanalbund, da DHM inkluderer befæstede overflader. Korrekt håndtering
        af sådanne krydsninger kræver supplerende data (fx ledningsregistre).

  \item \textbf{Netværksforbindelse:} Trin 2 kræver adgang til
        \texttt{api.dataforsyningen.dk} port 443. I netværksmiljøer med
        proxy eller VPN skal GDAL/QGIS konfigureres med korrekte
        proxyindstillinger.

  \item \textbf{DHM-opløsning:} Standard 2 m er passende for
        vådområdemodellering i hektar-skala. For meget store projektområder
        ($>50~\text{km}^2$) anbefales 5–10 m for at reducere procestiden.
\end{itemize}
